% Rúbrica: Introducción (abstract claro y conciso, en lenguaje accesible para un público general).
\section*{Resumen}
\addcontentsline{toc}{section}{Resumen}

\preguntaguia{} Esa es la pregunta que NexPlay ayuda a contestar. Steam deja devolver un videojuego
si se ha jugado menos de dos horas, y NexPlay usa ese límite: cuenta, en cada juego, las reseñas
negativas escritas antes de las dos horas, una señal de que el juego decepcionó pronto. Con \cifraResenasEntrenamiento{} reseñas de
\cifraJuegosEntrenamiento{} juegos se entrenó un modelo sencillo que estima qué juegos dejan más
de esa señal usando solo lo que cualquiera ve en la tienda antes de pagar: si el juego es gratis,
su precio, su descuento y la nota de la crítica. El modelo distingue a esos juegos mejor que un
clasificador que no sabe nada, también en \cifraJuegosExternos{} juegos que no vio, aunque la
ventaja es modesta. El sitio no da probabilidades ni dice qué comprar: da un riesgo de
arrepentimiento bajo, medio o alto; explica qué lo mueve y de qué se quejan las reseñas; deja
comparar juegos y preguntarle a un asistente, y recuerda la ventana de reembolso. Es una segunda
opinión antes de pagar. Un análisis aparte, con modelos que leen las reseñas, mostró que esas quejas tempranas
tienen un lenguaje propio: hablan de entrar al juego, del rendimiento y del reembolso. El riesgo,
sin embargo, sale solo de datos del juego y es el mismo para cualquier persona.
